\chapter{The mean--covariance term}\label{sec:th}

In this section, we bound the mean--covariance term
$J=\int j_q\,\nu(\dd q)$ at a minimizing profile. Here
$j_q=N^{-2}\E[m_q^TC_qm_q]$ is the expected conditional variance of the
spin in the direction of its conditional mean. Our main result,
\cref{prop:th-j}, states that $J\le341(1+\log_+M)^2/(Ms)$ at every
minimizer. Together with the trace bound $A\le1/(Ms)$ of
\eqref{eq:in-Lvar}, this controls the thermal part $A+2J$ of the defect
$\calD_N(\rho)=\beta(A+2J+V)$. As
$Ms=\beta^2N\rho\cdot\lambda/\beta=\beta N\rho\lambda$, the two bounds
imply that
\[
 \beta(A+2J)\le\frac{1+682(1+\log_+M)^2}{N\rho\lambda}.
\]
We use \cref{prop:th-j} in this form in the proof of
\cref{cor:c7-defect}, for those $\rho$ where the cavity inequality
contracts, and in \eqref{eq:sm-AJ} in the proof of \cref{prop:sm-small},
for small $\rho$.

The proof has four steps.
\begin{enumerate}
\item In \cref{lem:th-ric}, we show by a scalar Riccati inequality that
conditional covariances decay in every direction that is measurable at
an earlier level: if $b$ is $\calG_{q'}$-measurable, then
$\E[b^TC_qb]$ is at most $\E|b|^2$ divided by the accumulated clock
$c\int_{q'}^qa(u)^2\dd u$. This holds for every admissible profile.
\item In \cref{lem:th-gamma}, we use first-order optimality to show that
$\mathsf S(r)=\E S_r$ is $2$-Lipschitz toward every point $q$ of
$\supp\nu$, also when $r$ lies in a gap of the support. By
\eqref{eq:th-orth} below, this bounds $\E|m_q-m_r|^2$ by $2N|q-r|$.
\item In \cref{lem:th-res}, we combine these two estimates through a
dyadic decomposition of $m_q-m_p$, for $q\in\supp\nu$ and $p\le q$. We
bound $\calR(p,q)=N^{-2}\E[(m_q-m_p)^TC_q(m_q-m_p)]$ by a multiple of
$(1+\log_+M)^2/(Ma((p+q)/2)^2)$.
\item In \cref{prop:th-j}, we average earlier means $m_{Q(w)}$, $0<w<v$,
over the whole quantile prefix with a quadratic weight, together with a
deterministic weight on $m_0=0$. We control the conditional covariance in the direction of this average by
\eqref{eq:in-Lvar}, and
its distance from $m_{Q(v)}$ by \cref{lem:th-res}. The weight produces a factor $(s+v)^{-2}$, whose
integral over $v\in(0,1-s)$ is at most $1/s$.
\end{enumerate}

We use \cref{prop:fnd-contact}, namely the martingale equation and
covariance order \eqref{eq:in-martingale}, the first variation
\eqref{eq:fnd-first-variation}, and the contact identities
\eqref{eq:in-contact}; \cref{prop:fnd-variance}; and facts about the
quantile function from \cref{sec:wi-setup}, including
\eqref{eq:wi-supp}. \Cref{lem:th-ric,lem:th-gamma,lem:th-res} are used
only in this section.

\section{Standing assumptions}\label{sec:th-setup}

Throughout this section, we fix $N\ge1$, $0<\rho\le1$, and $0<s\le1/2$,
so that $t=1-\rho\in[0,1)$, $c>0$, and $M=Nc>0$. Expectations refer to the joint law of the
deterministic admissible profile $\mu=s\delta_0+\nu$ under
consideration. \Cref{lem:th-ric} holds for every admissible profile.
The later results assume that $\mu$ minimizes $\mathcal F_{t,s}$; they
hold for every minimizer, with absolute constants.

We write $\mathsf S(q)=\E S_q$ for the expected self-overlap. Recall
from \eqref{eq:in-martingale} that $a(q)C_q\preceq H_q\preceq C_q$, and
in particular $0\preceq H_q\preceq C_q$. The diagonal entries of $C_q$
are $1-m_{q,i}^2\le1$. These two facts give
\begin{equation}\label{eq:th-op}
 \|H_q\|_{\rm op}\le\|C_q\|_{\rm op}\le\Tr C_q\le N,\qquad
 0\le\mathsf S(q)\le1,
\end{equation}
where $\|\cdot\|_{\rm op}$ is the operator norm. The process $m$ is a
bounded martingale, and $\E[m_q\mid\calG_p]=m_p$ for $p\le q$. Then
$\E|m_q|^2=\E|m_p|^2+\E|m_q-m_p|^2$. In particular, $\mathsf S$ is
nondecreasing, and
\begin{equation}\label{eq:th-orth}
 \E|m_q-m_p|^2=N\bigl(\mathsf S(q)-\mathsf S(p)\bigr)
 \qquad(0\le p\le q\le1).
\end{equation}
For $0\le p\le q\le1$, we define
\[
 \calR(p,q)=N^{-2}\E\bigl[(m_q-m_p)^TC_q(m_q-m_p)\bigr],
\]
which is the analog of $j_q$ for the increment $m_q-m_p$ in place of
$m_q$. Note that $\calR(0,q)=j_q$, since $m_0=0$ by
\eqref{eq:in-martingale}.

\section{Covariance along an earlier measurable direction}\label{sec:th-decay}

Fix a direction that is measurable at some level. As more of the
auxiliary field is revealed, the expected conditional variance of the
spin in this direction can only decrease. The martingale equation
determines the rate of decrease, and the covariance order bounds this
rate from below by a multiple of the square of the variance itself. We integrate
the resulting scalar Riccati inequality and obtain decay that is
inversely proportional to the accumulated clock.

\begin{lemma}[Directional covariance decay]\label{lem:th-ric}
Let $\mu=s\delta_0+\nu$ be any deterministic admissible profile, let
$0\le q'<q\le1$, and let
\[
 \mathcal I=c\int_{q'}^qa(u)^2\dd u .
\]
Then $\mathcal I\ge cs^2(q-q')>0$, and for every
$\calG_{q'}$-measurable random vector $b\in\R^N$ with $\E|b|^2<\infty$,
\[
 \E[b^TC_qb]\le\frac{\E|b|^2}{\mathcal I}.
\]
\end{lemma}

\begin{proof}
The lower bound for $\mathcal I$ holds because $a\ge s$. We first prove
the main bound for bounded $b$, by deriving a differential inequality for
the function $r\mapsto\E[b^TC_rb]$ on $[q',q]$, and we then remove the
boundedness by monotone convergence.

\paragraph{Step 1: the rate of decrease.}
Suppose that $b$ is bounded, and define $\varphi(r)=\E[b^TC_rb]$ for
$r\in[q',q]$. For such $r$, the vector $b$ is $\calG_r$-measurable, so
that $b^TC_rb=\E[(b\cdot(\sigma-m_r))^2\mid\calG_r]$ and
$\varphi(r)=\E(b\cdot(\sigma-m_r))^2$. We write
\[
 b\cdot(\sigma-m_{q'})=b\cdot(\sigma-m_r)+b\cdot(m_r-m_{q'}).
\]
The second term is $\calG_r$-measurable, and the first has conditional
mean zero given $\calG_r$. The cross term has expectation
zero, and
\[
 \varphi(q')=\varphi(r)+\E\bigl(b\cdot(m_r-m_{q'})\bigr)^2 .
\]
By \eqref{eq:in-martingale}, and since $b$ is $\calG_{q'}$-measurable
and bounded, we have
$b\cdot(m_r-m_{q'})=\sqrt c\int_{q'}^r(H_ub)\cdot\dd B_u$. The integrand
is progressive and bounded by $N|b|$, by \eqref{eq:th-op}. Itô's isometry
gives
\[
 \varphi(r)=\varphi(q')-c\int_{q'}^r\E|H_ub|^2\dd u\qquad(q'\le r\le q).
\]
Then $\varphi$ is absolutely continuous and nonincreasing, with
$\varphi'(r)=-c\E|H_rb|^2$ for almost every $r\in[q',q]$.

\paragraph{Step 2: the Riccati inequality.}
On the event $\{b\ne0\}$, the Cauchy--Schwarz inequality and the
covariance order $0\le a(r)b^TC_rb\le b^TH_rb$ imply that
\[
 |H_rb|^2\ge\frac{(b^TH_rb)^2}{|b|^2}\ge\frac{a(r)^2(b^TC_rb)^2}{|b|^2}.
\]
Let $Z_r=|b|^{-1}b^TC_rb$ on $\{b\ne0\}$ and $Z_r=0$ otherwise. Then
$\E|H_rb|^2\ge a(r)^2\E Z_r^2$. Further, $\varphi(r)=\E[Z_r|b|]$, and by the
Cauchy--Schwarz inequality, $\varphi(r)^2\le\E Z_r^2\,\E|b|^2$. If
$\E|b|^2>0$, we conclude that
\[
 \varphi'(r)\le-\frac{c\,a(r)^2\varphi(r)^2}{\E|b|^2}
 \qquad\text{for almost every }r\in[q',q].
\]

\paragraph{Step 3: integration.}
If $\E|b|^2=0$ or $\varphi(q)=0$, there is nothing to prove. Otherwise,
$\varphi\ge\varphi(q)>0$ on $[q',q]$, because $\varphi$ is nonincreasing.
The function $1/\varphi$ is then absolutely continuous, with
$(1/\varphi)'=-\varphi'/\varphi^2\ge c\,a(r)^2/\E|b|^2$ almost
everywhere. Integrating over $[q',q]$, we find
\[
 \frac1{\varphi(q)}\ge\frac1{\varphi(q')}+\frac{\mathcal I}{\E|b|^2}
 \ge\frac{\mathcal I}{\E|b|^2},
\]
which is the claim for bounded $b$. The profile $a$ may have jumps, but
only Lebesgue-almost every level enters this integration.

\paragraph{Step 4: general directions.}
For general $b$ with $\E|b|^2<\infty$, let $b_n=b\ind_{\{|b|\le n\}}$,
which is bounded and $\calG_{q'}$-measurable. The quantity
$b_n^TC_qb_n=(b^TC_qb)\ind_{\{|b|\le n\}}$ increases to $b^TC_qb$ as
$n\to\infty$, because $C_q\succeq0$, and $|b_n|^2$ increases to
$|b|^2$. The bound holds for each $b_n$, and monotone convergence as
$n\to\infty$ yields the claim for $b$.
\end{proof}

\section{Mean increments toward a support point}\label{sec:th-support}

In the dyadic argument below, we need $\E|m_q-m_r|^2$ to be at most of
order $N|q-r|$ when $q$ is a point of $\supp\nu$, including when $r$
lies in a gap of the support. In view of \eqref{eq:th-orth}, this is a
Lipschitz bound for $\mathsf S$ toward $q$, and it follows from the global
first-order optimality of a minimizer. Let
\[
 G(y)=\int_y^1(\mathsf S(u)-u)\dd u\qquad(0\le y\le1)
\]
be the function in the first variation \eqref{eq:fnd-first-variation}.
\Cref{prop:fnd-contact} states that if $\mu$ is a minimizer, then
$\supp\nu$ is contained in the set of minimum points of $G$ on $[0,1]$,
that is,
\begin{equation}\label{eq:th-argmin}
 G(y)\ge G(q)\qquad(q\in\supp\nu,\ 0\le y\le1).
\end{equation}
We recall the argument from the proof of \cite[Proposition~2.1]{AL},
because we will use \eqref{eq:th-argmin} at points $y$ outside
$\supp\nu$. For $y\in[0,1]$, the profiles
$s\delta_0+(1-\theta)\nu+\theta(1-s)\delta_y$, $0\le\theta\le1$, are
admissible. At a minimizer, the right derivative of $\mathcal F_{t,s}$
along them at $\theta=0$, which is \eqref{eq:fnd-first-variation} with
$\nu'=(1-s)\delta_y$, is nonnegative. Dividing by $c/2$, we
find that
\[
 (1-s)G(y)\ge\int G(p)\,\nu(\dd p)\qquad(0\le y\le1).
\]
The function $G$ is continuous. Taking $y$ to be a minimum point of $G$
and recalling that $\nu$ has mass $1-s$, we find that
$\int(G-\min G)\dd\nu\le0$. Then the open set $\{G>\min G\}$ is
$\nu$-null and therefore does not meet $\supp\nu$.

\begin{lemma}[Lipschitz bound toward the support]\label{lem:th-gamma}
Let $\mu$ be a minimizer of $\mathcal F_{t,s}$. For every $q\in\supp\nu$
and every $r\in[0,1]$,
\[
 |\mathsf S(r)-\mathsf S(q)|\le2|r-q|;\qquad\text{in particular,}\qquad
 0\le\mathsf S(q)-\mathsf S(r)\le2(q-r)\quad(0\le r\le q).
\]
\end{lemma}

\begin{proof}
Let $q\in\supp\nu$, so that $\mathsf S(q)=q$ by \eqref{eq:in-contact}.
Recall that $\mathsf S$ is nondecreasing with values in $[0,1]$, and
note that $G(y_1)-G(y_2)=\int_{y_1}^{y_2}(\mathsf S(u)-u)\dd u$ for
$y_1\le y_2$. Let $r\in[0,1]$ and $\xi=\mathsf S(r)\in[0,1]$. We compare
the values of $G$ at $\xi$ and at $q$. The point $\xi$ need not lie in
$\supp\nu$, and this is why we need \eqref{eq:th-argmin} for every
$y\in[0,1]$.

\paragraph{The case $r\le q$.}
By monotonicity, $\mathsf S(q)-\mathsf S(r)\ge0$. If $\xi\ge r$, then
$\mathsf S(q)-\mathsf S(r)=q-\xi\le q-r$. If $\xi<r$, then
$\mathsf S\le\mathsf S(r)=\xi$ on $[\xi,r]$ and
$\mathsf S\le\mathsf S(q)=q$ on $[r,q]$, by monotonicity. By
\eqref{eq:th-argmin} with $y=\xi$, we have
\[
 \begin{aligned}
 0\le G(\xi)-G(q)
 &=\int_\xi^r(\mathsf S(u)-u)\dd u+\int_r^q(\mathsf S(u)-u)\dd u\\
 &\le\int_\xi^r(\xi-u)\dd u+\int_r^q(q-u)\dd u
 =\frac{(q-r)^2-(r-\xi)^2}2 .
 \end{aligned}
\]
Then $r-\xi\le q-r$, and
$\mathsf S(q)-\mathsf S(r)=(q-r)+(r-\xi)\le2(q-r)$.

\paragraph{The case $r>q$.}
By monotonicity, $\mathsf S(r)-\mathsf S(q)\ge0$. If $\xi\le r$, then
$\mathsf S(r)-\mathsf S(q)=\xi-q\le r-q$. If $\xi>r$, then
$\mathsf S\ge q$ on $[q,r]$ and $\mathsf S\ge\xi$ on $[r,\xi]$, and in
the same way
\[
 \begin{aligned}
 0\le G(\xi)-G(q)
 &=-\int_q^r(\mathsf S(u)-u)\dd u-\int_r^\xi(\mathsf S(u)-u)\dd u\\
 &\le\int_q^r(u-q)\dd u+\int_r^\xi(u-\xi)\dd u
 =\frac{(r-q)^2-(\xi-r)^2}2 .
 \end{aligned}
\]
We obtain $\xi-r\le r-q$, and
$\mathsf S(r)-\mathsf S(q)=(r-q)+(\xi-r)\le2(r-q)$. This completes the
proof.
\end{proof}

Together, \eqref{eq:th-orth} and \cref{lem:th-gamma} imply that
\[
 \E|m_q-m_r|^2\le2N(q-r)\qquad(q\in\supp\nu,\ 0\le r\le q).
\]

\section{A dyadic decomposition of the mean increment}\label{sec:th-dyadic}

\Cref{lem:th-ric} does not apply directly to the increment $m_q-m_p$,
because this increment is not measurable at any level before $q$.
Instead, we split it into increments whose right endpoints approach $q$
geometrically. Each of them except the last is measurable strictly
before $q$, and \cref{lem:th-ric} applies to it. By \cref{lem:th-gamma},
its second moment is at most proportional to the remaining distance to $q$,
which cancels the length of the clock in \cref{lem:th-ric}. The last,
short increment is controlled by $C_q\preceq NI$.

\begin{lemma}[Covariance of a mean increment]\label{lem:th-res}
Let $\mu$ be a minimizer of $\mathcal F_{t,s}$, let $q\in\supp\nu$ and
$0\le p\le q$, and let $a_*=a((p+q)/2)$. For every integer $K\ge1$,
\begin{equation}\label{eq:th-res-K}
 \calR(p,q)\le(K+1)\Bigl[\frac{4K}{Ma_*^2}+2(q-p)2^{-K}\Bigr].
\end{equation}
Moreover,
\begin{equation}\label{eq:th-res}
 \calR(p,q)\le\frac{56\bigl(1+\log_+(M(q-p))\bigr)^2}{Ma_*^2}.
\end{equation}
\end{lemma}

\begin{proof}
We may assume that $p<q$, since for $p=q$ the left side vanishes and
both bounds hold. The proof has three steps. We split $m_q-m_p$ into
$K+1$ dyadic pieces; we bound the first $K$ pieces by
\cref{lem:th-ric,lem:th-gamma} and the last piece directly; and we then
choose $K$.

For $0\le k\le K$, let $q_k=q-(q-p)2^{-k}$, so that $q_0=p$,
$q_1=(p+q)/2$, and $q-q_k=(q-p)2^{-k}$. We set
$b_k=m_{q_{k+1}}-m_{q_k}$ for $0\le k<K$ and $b_K=m_q-m_{q_K}$. Then
$m_q-m_p=\sum_{k=0}^Kb_k$. Applying the Cauchy--Schwarz inequality to the
quadratic form of $C_q$, which is positive semidefinite, we find
\[
 (m_q-m_p)^TC_q(m_q-m_p)\le(K+1)\sum_{k=0}^Kb_k^TC_qb_k .
\]
The points $q_k$ need not be contact levels; we use \cref{lem:th-gamma}
only with the endpoint $q\in\supp\nu$.

\paragraph{Step 1: the first $K$ pieces.}
Let $0\le k<K$. The vector $b_k$ is $\calG_{q_{k+1}}$-measurable, and
$q_{k+1}<q$. Since $a$ is nondecreasing and $q_{k+1}\ge q_1$, we have
$a(u)\ge a(q_1)=a_*$ for $u\ge q_{k+1}$. The clock of \cref{lem:th-ric}
with $q'=q_{k+1}$ satisfies
\[
 \mathcal I=c\int_{q_{k+1}}^qa(u)^2\dd u\ge ca_*^2(q-q_{k+1}),
\]
and \cref{lem:th-ric} implies that
$\E[b_k^TC_qb_k]\le\E|b_k|^2/(ca_*^2(q-q_{k+1}))$. By \eqref{eq:th-orth},
the monotonicity of $\mathsf S$, and \cref{lem:th-gamma} with $r=q_k$, we
have
\[
 \E|b_k|^2=N\bigl(\mathsf S(q_{k+1})-\mathsf S(q_k)\bigr)
 \le N\bigl(\mathsf S(q)-\mathsf S(q_k)\bigr)\le2N(q-q_k)=4N(q-q_{k+1}).
\]
The last two bounds imply that
\[
 N^{-2}\E[b_k^TC_qb_k]\le\frac4{Nca_*^2}=\frac4{Ma_*^2}
 \qquad(0\le k<K).
\]

\paragraph{Step 2: the last piece.}
Using $\|C_q\|_{\rm op}\le N$ from \eqref{eq:th-op}, then
\eqref{eq:th-orth} and \cref{lem:th-gamma} with $r=q_K$, we have
\[
 N^{-2}\E[b_K^TC_qb_K]\le N^{-1}\E|b_K|^2=\mathsf S(q)-\mathsf S(q_K)
 \le2(q-q_K)=2(q-p)2^{-K}.
\]
We sum the bounds of Steps~1 and~2 over the $K+1$ pieces and insert
them into the Cauchy--Schwarz bound above. This proves
\eqref{eq:th-res-K}.

\paragraph{Step 3: the choice of $K$.}
Let $z=M(q-p)$, $\theta=\log_+z$, and $K=\lceil\log_2(2+z)\rceil$, so that
$K\ge1$. Then $2^K\ge2+z>z$, which gives
\[
 2(q-p)2^{-K}=\frac{2z}{M2^K}\le\frac2M\le\frac2{Ma_*^2},
\]
where we used $a_*\le1$. With this choice of $K$, \eqref{eq:th-res-K}
implies that
\[
 \calR(p,q)\le\frac{(K+1)(4K+2)}{Ma_*^2}.
\]
Since $2+z\le3\max\{1,z\}$, we have
$\log_2(2+z)\le\log_23+\theta/\log2<2+2\theta$, where we used
$\log_23<2$ and $1/\log2<2$. It follows that
$K<1+\log_2(2+z)<3+2\theta$, and
\[
 (K+1)(4K+2)\le(4+2\theta)(14+8\theta)=56+60\theta+16\theta^2
 \le56(1+\theta)^2 .
\]
This proves \eqref{eq:th-res}.
\end{proof}

The two powers of the logarithm in \eqref{eq:th-res} have different
sources: there are $K$ increments of comparable cost, and the
Cauchy--Schwarz inequality for their sum contributes the factor $K+1$.

\section{An integrated bound for \texorpdfstring{$j$}{j}}\label{sec:th-integrated}

With $p=0$, \cref{lem:th-res} bounds $j_q=\calR(0,q)$ by a multiple of
$(1+\log_+M)^2/(Ma(q/2)^2)$. To integrate against $\nu$, we need a bound
in terms of the cumulative mass $s+v$ at $q=Q(v)$, rather than
$a(q/2)$, which may be as small as $s$. We therefore average the earlier
means $m_{Q(w)}$, $0<w<v$, with density proportional to $(s+w)^2$, and
assign the remaining weight to $m_0=0$. The conditional covariance in the direction of
this average is controlled by the integrated variance bound
\eqref{eq:in-Lvar}, and its distance from $m_{Q(v)}$ by
\cref{lem:th-res}. With the quadratic weight, both estimates are of
order $(s+v)^{-2}$, also for $v<s$, and no separate argument is needed
near the root.

\begin{proposition}[Covariance in the mean direction]\label{prop:th-j}
Let $\mu=s\delta_0+\nu$ be a minimizer of $\mathcal F_{t,s}$. For every
$v\in(0,1-s)$,
\begin{equation}\label{eq:th-j-pointwise}
 j_{Q(v)}\le\frac{341(1+\log_+M)^2}{M(s+v)^2},
\end{equation}
and
\begin{equation}\label{eq:th-J}
 J=\int j_q\,\nu(\dd q)\le\frac{341(1+\log_+M)^2}{Ms}.
\end{equation}
The constant is absolute. The bounds hold for all $N\ge1$, $c>0$,
$0<s\le1/2$, and every minimizing profile, including profiles with atoms.
\end{proposition}

\begin{proof}
Fix $v\in(0,1-s)$, and let $q=Q(v)$ and $\mathcal L=1+\log_+M$. By
\eqref{eq:wi-supp}, we have $q\in\supp\nu$, and $Q(w)\le q$ for $w<v$
because $Q$ is nondecreasing. If $q=0$, then $m_q=m_0=0$ and $j_q=0$, so
we may assume that $q>0$. The proof has four steps. First, we write
$m_q$ as a weighted average $\bar m$ of earlier means plus a remainder
$\bar\Delta$, and we prove the bound \eqref{eq:th-split} below. Second,
we bound the conditional covariance in the direction of $\bar m$ by
\eqref{eq:in-Lvar},
using the conditional Jensen inequality on quantile fibers. Third, we
bound the remainder by \cref{lem:th-res}. Fourth, we combine the bounds
and integrate over $v$.

\paragraph{Step 1: the decomposition.}
We use the deterministic weights
\[
 r_0=\frac{s^3}{(s+v)^3},\qquad
 k_v(w)=\frac{3(s+w)^2}{(s+v)^3}\quad(0<w<v),
\]
which satisfy
\[
 r_0+\int_0^vk_v(w)\dd w=\frac{s^3+\bigl((s+v)^3-s^3\bigr)}{(s+v)^3}=1 .
\]
We use the probability measure $r_0\delta_0+k_v(w)\dd w$ on $[0,v)$ to
average the vector $m_0=0$ and the vectors $m_{Q(w)}$, $0<w<v$. We define
\[
 \begin{gathered}
 \bar m=\int_0^vk_v(w)m_{Q(w)}\dd w,\qquad
 \bar L=\int_0^vk_v(w)L_{Q(w)}\dd w,\\
 \bar\Delta=r_0m_q+\int_0^vk_v(w)\bigl(m_q-m_{Q(w)}\bigr)\dd w.
 \end{gathered}
\]
By the normalization of the weights, $m_q=\bar m+\bar\Delta$. The vector
$\bar m$ and the number $D=2^{-1}\int_0^vk_v(w)|m_{Q(w)}|^2\dd w$ are
$\calG_q$-measurable, because $Q(w)\le q$ for $w<v$, and
$N\bar L=\sigma\cdot\bar m-D$. It follows that
$\Var(\bar L\mid\calG_q)=N^{-2}\bar m^TC_q\bar m$, and by the law of
total variance,
\[
 N^{-2}\E[\bar m^TC_q\bar m]=\E\Var(\bar L\mid\calG_q)\le\Var(\bar L).
\]
Next, we write
$\bar\Delta=r_0(m_q-m_0)+\int_0^vk_v(w)(m_q-m_{Q(w)})\dd w$, which is
the average of the vectors $m_q-m_0$ and $m_q-m_{Q(w)}$ with respect to
the probability measure above. By Jensen's inequality for the
convex quadratic form of $C_q$, we have
\[
 \bar\Delta^TC_q\bar\Delta
 \le r_0(m_q-m_0)^TC_q(m_q-m_0)
 +\int_0^vk_v(w)(m_q-m_{Q(w)})^TC_q(m_q-m_{Q(w)})\dd w .
\]
The integrand is bounded by $4N^2$, by \eqref{eq:th-op}. We take
expectations, use Fubini's theorem, and divide by $N^2$. This gives
\[
 N^{-2}\E[\bar\Delta^TC_q\bar\Delta]
 \le r_0\calR(0,q)+\int_0^vk_v(w)\calR(Q(w),q)\dd w .
\]
Finally, since $C_q\succeq0$, we have
$m_q^TC_qm_q\le2\bar m^TC_q\bar m+2\bar\Delta^TC_q\bar\Delta$. Combining
the last three bounds, we obtain
\begin{equation}\label{eq:th-split}
 j_q\le2\Var(\bar L)
       +2r_0\calR(0,q)+2\int_0^v k_v(w)\calR(Q(w),q)\dd w.
\end{equation}
The quantity $\calR(0,q)=j_q$ appears on the right side. This is not
circular: in Step~3, we bound it by \cref{lem:th-res}, independently of
\eqref{eq:th-split}.

\paragraph{Step 2: the variance of the weighted prefix.}
We extend $k_v$ by zero to $(0,1-s)$, and note that $k_v\le3/(s+v)$.
The image of the measure $k_v(w)\dd w$ under $Q$ is then at most
$3/(s+v)$ times $\nu$, and it has a density $\kappa_v$ with respect to
$\nu$, with $0\le\kappa_v\le3/(s+v)$; in particular,
$\kappa_v\in L^2(\nu)$. By the definition of $\kappa_v$ and the pushforward property of $Q$, we have,
for every bounded Borel function $\varphi$,
\[
 \int_0^{1-s}k_v(w)\varphi(Q(w))\dd w=\int\kappa_v\varphi\dd\nu
 =\int_0^{1-s}\kappa_v(Q(w))\varphi(Q(w))\dd w .
\]
Applied for each outcome to the bounded function $p\mapsto L_p$, the
first identity shows that $\bar L=\int\kappa_v(p)L_p\,\nu(\dd p)$. The identity as a whole says
that, under the normalized Lebesgue measure on $(0,1-s)$, the function
$\kappa_v\circ Q$ is the conditional expectation of $k_v$ given the
$\sigma$-algebra generated by $Q$. Explicitly, if $p$ is an atom of
$\nu$ with mass $r_p$, then
\[
 \kappa_v(p)=\frac1{r_p}\int_{(F_-(p),F(p))}k_v(w)\dd w,
\]
with the notation of \cref{sec:wi-setup}. A prefix $(0,v)$ that ends
inside the quantile interval of an atom gives the atom the average
of $k_v$ over that interval. Since $1/a(Q)$ is a bounded nonnegative
function of $Q$, we may apply the conditional Jensen inequality and find
\[
 \int\frac{\kappa_v(p)^2}{a(p)}\,\nu(\dd p)
 =\int_0^{1-s}\frac{\kappa_v(Q(w))^2}{a(Q(w))}\dd w
 \le\int_0^{1-s}\frac{k_v(w)^2}{a(Q(w))}\dd w
 =\int_0^v\frac{k_v(w)^2}{a(Q(w))}\dd w .
\]
As $\mu$ is a minimizer, \cref{prop:fnd-variance} applies with
$k=\kappa_v$. Together with $a(Q(w))\ge s+w$, this yields
\[
 \begin{aligned}
 \Var(\bar L)
 &\le\frac1M\int_0^v\frac{k_v(w)^2}{s+w}\dd w
 =\frac9{M(s+v)^6}\int_0^v(s+w)^3\dd w\\
 &=\frac9{4M(s+v)^2}\Bigl[1-\frac{s^4}{(s+v)^4}\Bigr]
 \le\frac9{4M(s+v)^2}.
 \end{aligned}
\]

\paragraph{Step 3: the remainder.}
Let $0<w<v$ and $p=Q(w)$. Then $p\le q$, and since $a$ is
nondecreasing, $a((p+q)/2)\ge a(p)\ge s+w$. For the root term, we have
$a(q/2)\ge s$ by \eqref{eq:wi-supp}. All level differences are at most
one, and hence $\log_+(M(q-p))\le\log_+M$. Applying \cref{lem:th-res}
with the endpoint $q\in\supp\nu$ and the lower endpoints $0$ and $Q(w)$, we have
\[
 \calR(0,q)\le\frac{56\mathcal L^2}{Ms^2},\qquad
 \calR(Q(w),q)\le\frac{56\mathcal L^2}{M(s+w)^2}\quad(0<w<v).
\]
Using $r_0s^{-2}=s/(s+v)^3$ and $k_v(w)(s+w)^{-2}=3/(s+v)^3$, we find
that
\[
 r_0\calR(0,q)+\int_0^vk_v(w)\calR(Q(w),q)\dd w
 \le\frac{56\mathcal L^2}{M}\cdot\frac{s+3v}{(s+v)^3}
 \le\frac{168\mathcal L^2}{M(s+v)^2},
\]
where we used $s+3v\le3(s+v)$.

\paragraph{Step 4: conclusion.}
Inserting the bounds of Steps~2 and~3 into \eqref{eq:th-split} and using
$\mathcal L\ge1$, we obtain
\[
 j_{Q(v)}\le\frac9{2M(s+v)^2}+\frac{336\mathcal L^2}{M(s+v)^2}
 \le\frac{(9/2+336)\mathcal L^2}{M(s+v)^2}
 \le\frac{341\mathcal L^2}{M(s+v)^2},
\]
which is \eqref{eq:th-j-pointwise}. Finally,
$J=\int_0^{1-s}j_{Q(v)}\dd v$, because $Q$ pushes Lebesgue measure on
$(0,1-s)$ forward to $\nu$, and \eqref{eq:th-J} follows from
\[
 \int_0^{1-s}\frac{\dd v}{(s+v)^2}=\frac1s-1\le\frac1s .
\]
This completes the proof.
\end{proof}
